\section{The lookup tables}
\label{sec:tables}

The known algorithms live in JSON, one file per stage, inside the package:

\begin{center}
\begin{tabular}{@{}llr@{}l@{}}
\toprule
File & Stage & Cases & Labels \\
\midrule
\texttt{data/oll.json} & orientation of the last layer & 57 & \texttt{"1"} \ldots{} \texttt{"57"} \\
\texttt{data/pll.json} & permutation of the last layer & 21 & \texttt{"Aa Perm"} \ldots{} \texttt{"Z Perm"} \\
\texttt{data/f2l.json} & the four corner--edge pairs & 41 & \texttt{"1"} \ldots{} \texttt{"41"} \\
\bottomrule
\end{tabular}
\end{center}

\subsection{Schema}

All three files share one shape:

\begin{lstlisting}[style=data]
{
  "schema": 1,
  "kind": "oll",
  "cases": [
    {
      "label": "1",
      "name": "No Edges - Runway",
      "algorithm": "R U2' R2' F R F' U2 R' F R F'"
    }
  ]
}
\end{lstlisting}

\begin{description}[leftmargin=0em,style=nextline]
  \item[\texttt{schema}] the format version, so a future change can be detected
  rather than misread;
  \item[\texttt{kind}] \texttt{oll}, \texttt{pll} or \texttt{f2l}, which tells the
  loader which recognizer to build;
  \item[\texttt{label}] what the user sees in the solution listing --- the OLL
  number, the PLL name, or the F2L case number;
  \item[\texttt{name}] a longer description, shown by interfaces and in the JSON
  metadata;
  \item[\texttt{algorithm}] one algorithm in WCA notation.
\end{description}

\subsection{Where the algorithms come from}

The tables are generated, not typed by hand, from the published \emph{Logiqx
cubing-algs} collection, whose case identifiers match the community numbering:

\begin{lstlisting}[style=shell]
uv run python tools/build_last_layer_tables.py oll.js pll.js f2l.js
\end{lstlisting}

For every case the tool considers all the algorithms listed for it, keeps the ones
that parse, checks that each one really solves its own case, and prefers the
shortest algorithm that uses face turns only.  Those preferences matter: a wide or
slice move shifts the centers, so an algorithm that uses one would leave the cube
in a different orientation than the net is drawn in.

\subsection{Why recognition cannot drift}

The loader derives each case's recognition key from its algorithm, by applying the
inverse of the algorithm to a solved cube and reading the pattern of the result.  A
consequence worth knowing: if you edit an \texttt{algorithm} and restart, both the
solution and the recognition change together, and there is no second place to update.

Two cases whose patterns collide are a hard error --- the loader raises
\texttt{InvalidStateError} naming both labels --- because a state with two possible
algorithms is ambiguous.

\subsection{Adding or changing a case}

\begin{enumerate}[leftmargin=1.4em]
  \item edit the JSON file: add an entry with a fresh label, a name and an algorithm;
  \item run the tests.  The per-case test applies every algorithm to its own case
  state and requires the stage to reach its goal, so a wrong algorithm fails loudly;
  \item run \texttt{scripts/check.sh} and commit.
\end{enumerate}

Tables are loaded once per process and cached, so a running program keeps the tables
it started with --- restart to pick up an edit.

\subsection{What the tables do not cover}

Two gaps are worth stating plainly:

\begin{itemize}[leftmargin=1.4em]
  \item of the 41 published F2L cases, 35 survive the ``face turns only''
  preference.  Nothing is lost: the other six, and every configuration the 41 cases
  do not describe at all, are handled by the verified search.
  \item the last layer can contain a state for which the stored algorithm does not
  finish the cube.  The pipeline notices, finishes the remainder with the two-phase
  solver and reports it honestly in the stage metadata.
\end{itemize}
